\section{基点を持つ四二次式系の支配的チャート}
\label{sec:product-net-boundary}
ここでは最初の三二次式も動かし、基点を持つ四次元二次式空間の一般点を覆う。
この被覆はGrassmann多様体上の接写像で証明し、有限個の標本から正則性を推測しない。

$p=(0,0,0,1)$で消える二次式の九次元空間を$T_p$、その四次元部分空間の多様体を
$Y=\operatorname{Gr}(4,T_p)$とする。$\dim Y=20$である。
次の八係数族を取る：
\begin{align*}
 f_0&=e_0^2-e_0e_3-\alpha e_0e_1,&
 f_1&=e_1^2-e_1e_3-\beta e_1e_2,\\
 f_2&=e_2^2-e_2e_3-\gamma e_2e_0,&
 q&=e_0e_3+u e_1e_3+v e_2e_3+a e_0e_1+b e_0e_2+c e_1e_2.
\end{align*}
$H\subset\operatorname{GL}(V^*)$を、$p$で消える一次式の三次元空間を保つ部分群とする。
これは13次元の滑らかで既約な群であり、共通スカラーは$Y$上自明に作用する。

\begin{theorem}[八係数積型netチャートの支配性]\label{thm:product-net-dominance}
射$H\times\mathbb A^8\to Y$、$(g,\theta)\mapsto g\langle f_0,f_1,f_2,q\rangle$は支配的である。
さらに最初の三式が相異なる八点完全交叉で、第四式が$p$だけで消える部分へ制限しても支配的である。
これは二次式系の一般点の被覆であり、indexed AS同型類の次元やAS正則性の認定ではない。
\end{theorem}
\begin{proof}
四式は最初の三平方係数と$q$の非零係数から独立であり、全係数で$Y$への射が定まる。
係数値
\[
 (\alpha,\beta,\gamma,u,v,a,b,c)=(2,3,5,2,3,5,7,11)
\]
と$g=I$で接写像を計算する。
$T_p/L$の五次元商へ各四式の変分を射影すれば、
接空間$\operatorname{Hom}(L,T_p/L)$の20座標が得られる。
$H$の13個の行列単位変分と八係数の変分を並べた行列は$20\times21$であり、
列$0,1,\ldots,11,13,14,\ldots,20$による20次小行列式は
\[
 -\frac{696}{78125}\ne0
\]
となる。単項式順と商基底の選択、微分行列の全構成は
\path{computations/verify_product_net_dominance.py}に置く。
従って接写像は全射であり、滑らかな既約な$Y$への射は支配的である。

この点で最初の三式の八点は、$e_3=1$として
\[
\begin{gathered}
(0,0,0),\ (1,0,0),\ (0,1,0),\ (0,0,1),\\
(3,1,0),\ (1,0,6),\ (0,4,1),\ (-9/29,-19/29,-16/29)
\end{gathered}
\]
である。無限遠には、三変数の積の連鎖から
$\alpha\beta\gamma=1$でなければ非零解はなく、この点では$30\ne1$である。
上の八点は相異なり、従って三二次式の次数8の完全交叉を尽くす。
$q$の他の七点での値は
\[
1,\ 2,\ 3,\ 20,\ 61,\ 55,\ 2452/841
\]
と全て非零である。
この条件はパラメータの開条件なので、制限した射も同じ全射接写像を持ち、支配的である。
\end{proof}

\begin{proposition}[支配的チャートから全基点域への障害の移送]
\label{prop:product-net-universal-obstruction-transfer}
上記八変数の一般点における三次双対フレーム・テンソルの五次$d_2$階数を$r$とする。
$r<84$が厳密に証明されれば、基点を持つ$L$から作る三因子対称テンソルは
標準型Koszul AS点を一つも含まない。
より正確には、その有限Frobenius・四次正確部分で始点0の$H_2$次元は全て$84-r$以上である。
ここで$L$に基点が一つだけであることや、相異なる八点netを持つことを仮定する必要はない。
\end{proposition}
\begin{proof}
まず基点を$p$に固定する。
$Y$上の可換二次商$B$について$\dim B_3\geq4$であり、等号となる部分$Y_3$は空でない開部分である。
実際、次数3の関係挿入は16本であり、Boolean標本で等号を達成する。
$Y_3$上の三次双対は階数4のベクトル束であり、そのフレーム束$\mathcal F$は既約である。
$\mathcal F$から三因子対称テンソルへの写像は正則である。
その有限Frobenius・四次正確部分$\mathcal U$は、二対二階数の上界6から開であり、既知標本を含む。

前定理の支配的射を$Y_3$へ制限し、さらに戻り因子の$\operatorname{GL}_4$基底変更を許せば、
$\mathcal F$への支配的な射が得られる。
八係数族で選んだ三次座標正規化が定まる部分も空でない開部分である。
共通の三因子変換と戻り因子の基底変更は、五次微分の階数を保存する。
従って八変数一般点の階数$r$は$\mathcal U$の一般階数でもある。
全ての$\mathcal U$上で$\operatorname{rk}d_2\leq r$、$\operatorname{rk}d_3\leq36$となり、
$\dim H_2\geq120-r-36=84-r$を得る。

最後に、有限Frobeniusな三因子対称テンソルでは、対称な三連続因子の三次交叉が
$B_3^*$に一致し、次元は4である。従って必ず上の$Y_3$のフレームから来る。
基点が他の場所なら共通の射影基底変換で$p$へ移せる。
AS点は有限Frobenius・四次正確部分に入るので、$r<84$なら全て排除される。
\end{proof}

この条件付き移行は一般点の精密な階数検算にも使えるが、基点枝全体のAS排除は
定理\ref{thm:partial-symmetric-basepoint-exclusion}で、一般点計算を介さずに得られた。
